\section{总结与延伸}

\subsection{Harness 设计的演化论}

\begin{importantbox}{核心论断：空间在移动，而非缩小}
\textit{``The space of interesting harness combinations doesn't shrink as models improve. Instead, it moves, and the interesting work for AI engineers is to keep finding the next novel combination.''}

有趣的 harness 组合空间不会随着模型改进而缩小——它会\textbf{移动}。随着模型变强，某些 harness 组件变得不再必要（如 Opus 4.6 下的 sprint 分解），但新的可能性也随之打开（如让 generator 构建嵌入 AI agent 的应用）。AI 工程师的核心工作不是被动等待模型进步使 harness 过时，而是\textbf{主动寻找下一个有价值的组合}。
\end{importantbox}

这一论断具有重要的方法论意义。它意味着 harness design 不是一个``最终会被模型进步淘汰''的过渡性工程领域，而是一个与模型能力\textbf{共同演进}的持久性工程领域。

\subsection{三条实践建议}

作者从这项工作中提炼了三条可携带的实践建议：

\textbf{1. 实验驱动，阅读 trace。} 用你正在构建的目标模型（而非假想中的模型）在真实问题上实验。阅读模型的执行 trace，理解它在哪里成功、在哪里失败。基于观察而非猜测来调优性能。

\textbf{2. 按需分解，专门化 agent。} 对于更复杂的任务，有时将任务分解并对每个方面应用专门的 agent 可以带来 headroom。但不要默认就分解——先验证单 agent 确实做不到。

\textbf{3. 模型升级时重审 harness。} 每当新模型发布，系统性地检查 harness 中的每个组件：剥离不再承载性能的组件（减少复杂度和成本），添加新组件以达成之前不可能的能力（利用新的模型能力）。

\subsection{结构化反思}

从本文的实践中可以提炼出几个层次的认知：

\textbf{方法论层面}：Harness 的每个组件都是一个可测试的假设（``模型不能独立完成 X''）。当模型进步时，假设需要重新验证。``假设-验证-迭代''构成了 harness 开发的核心循环。

\textbf{架构层面}：从 GAN 到 multi-agent 流水线的迁移展示了一个通用模式——将人类软件工程的最佳实践（代码审查、QA、sprint planning）\textbf{编码为 agent 角色}，而非试图让单个 agent 同时承担所有角色。角色分离使得每个 agent 的行为更易于独立调优。

\textbf{工程实践层面}：
\begin{itemize}[leftmargin=1.5em]
    \item 评分标准本身就是有效的 prompt engineering——即使没有 evaluator 反馈，仅将评分标准写入 generator 的 prompt 就能提升 baseline
    \item 文件作为 agent 间通信媒介（file-based communication）是一种简单、可追溯且有效的设计
    \item Planner 应保持高层次以避免错误级联，sprint contract 负责弥合细节间隙
    \item QA agent 的校准是一个需要多轮人工审查的调优过程，不存在``一次配置，永久有效''的捷径
    \item Evaluator 的价值与任务难度和模型能力之间的关系是\textbf{动态}的
\end{itemize}

\begin{knowledgebox}{与上一篇 Harness Design 文章的关系}
本文是 Anthropic 关于 harness design 系列的延续。上一篇文章（\textit{Harness design for agentic coding}）建立了 initializer + coding agent + context reset 的基本框架。本文在此基础上：
\begin{itemize}[leftmargin=1.5em]
    \item 引入了 GAN 启发的 generator-evaluator 分离
    \item 将主观设计质量纳入可评分框架
    \item 从 Sonnet 4.5 演进到 Opus 4.5/4.6，逐步简化架构
    \item 强调了 harness 组件的``保质期''概念——它们会随模型进步而过时
\end{itemize}
两篇文章共同构成了一个完整的 harness design 方法论。
\end{knowledgebox}

\subsection{拓展阅读}

\begin{itemize}[leftmargin=1.5em]
    \item Anthropic: \textit{Building Effective Agents} --- harness 设计的基础原则，``find the simplest solution possible''
    \item Anthropic: \textit{Context Engineering} --- 上下文窗口管理策略，context anxiety 的深入讨论
    \item Anthropic: \textit{Harness Design for Agentic Coding} --- 本文的前序文章，initializer + coding agent + context reset 架构
    \item \textbf{Claude Agent SDK} 文档 --- 构建 multi-agent 系统的官方框架
    \item \textbf{Playwright MCP} --- 让 agent 与真实浏览器交互的工具，evaluator 的核心能力来源
    \item Goodfellow et al., \textit{Generative Adversarial Networks} (2014) --- GAN 的原始论文，generator-discriminator 对抗训练的理论基础
\end{itemize}

%% --- Fin del contenido --- %%

\end{document}
